\section{Transpiling}\label{sec:transpiling}

\subsection{Quantum Fidelity}\label{sec:quantum-fidelity}

\subsubsection{Why Quantum Computers Are Noisy}\label{sec:why-quantum-computers-are-noisy}

In all the previous chapters, we implicitly assumed that quantum gates exactly as their mathematical definitions. For example:
\begin{itemize}
    \item Hadamard gate:
    \begin{equation*}
        H\ket{0} = \ket{+} = \frac{1}{\sqrt{2}}(\ket{0} + \ket{1})
    \end{equation*}
    Behaves as a perfect beam splitter, splitting the input state into an equal superposition of $\ket{0}$ and $\ket{1}$.

    \item Pauli-X gate:
    \begin{equation*}
        X\ket{0} = \ket{1}
    \end{equation*}
    Behaves as a perfect NOT gate, flipping the state of the qubit from $\ket{0}$ to $\ket{1}$.

    \item CNOT gate:
    \begin{gather*}
        \text{CNOT}\ket{00} = \ket{00}, \quad \text{CNOT}\ket{01} = \ket{01} \\
        \quad \text{CNOT}\ket{10} = \ket{11}, \quad \text{CNOT}\ket{11} = \ket{10}
    \end{gather*}
    Behaves as a perfect Controlled-NOT gate, flipping the target qubit if the control qubit is in the state $\ket{1}$ and leaving it unchanged if the control qubit is in the state $\ket{0}$.
\end{itemize}
\textbf{Mathematically} these transformations \textbf{are exact}. However, in a real quantum computer, gates are implemented by physical hardware:
\begin{itemize}
    \item Microwave pulses (superconducting qubits);
    \item Laser pulses (trapped ions);
    \item Electromagnetic fields;
    \item Control electronics;
\end{itemize}
Because the \textbf{hardware is imperfect} (e.g., due to manufacturing defects, environmental noise, or limitations in control precision), the \textbf{resulting operation is only an approximation of the ideal gate}. As a consequence:
\begin{equation*}
    U_{\text{physical}} \ne U_{\text{ideal}} \quad \xrightarrow{\text{better}} \quad U_{\text{physical}} \approx U_{\text{ideal}}
\end{equation*}
This discrepancy leads to \textbf{errors} in the quantum computation, which can accumulate over multiple gates and degrade the performance of quantum algorithms.

\newpage

\begin{flushleft}
    \textcolor{Red2}{\faIcon{exclamation-triangle} \textbf{The Two Main Sources of Noise}}
\end{flushleft}
\begin{itemize}
    \item \definition{Gate Infidelity}. Gate infidelity means that the \hl{gate physically applied by the hardware is not exactly the gate we intended to execute}.

    Suppose we want to apply a Hadamard gate:
    \begin{equation*}
        U_{\text{ideal}} = H
    \end{equation*}
    Ideally, we would like the physical implementation to match the ideal gate:
    \begin{equation*}
        H\ket{0} = \ket{+} = \frac{1}{\sqrt{2}}(\ket{0} + \ket{1})
    \end{equation*}
    However, due to \textbf{imperfections} like inaccurate control pulses, calibration errors, electronic noise, manufacturing imperfections, the hardware may actually apply:
    \begin{equation*}
        U_{\text{physical}} = H + \varepsilon
    \end{equation*}
    Where $\varepsilon$ is a \textbf{small error}. The resulting state is therefore slightly different from the desired one. After thousands of gates, these small errors can \textbf{accumulate} and significantly affect the final result.

    \highspace
    \textcolor{Green3}{\faIcon{question-circle} \textbf{Why Multi-Qubit Gates Are Worse?}} Intuitively, multi-qubit gates (like CNOT) are \textbf{more complex to implement} than single-qubit gates (like Hadamard) because they require precise control over interactions between multiple qubits. For this reason, the \textbf{fidelity of multi-qubit gates is typically lower} than that of single-qubit gates:
    \begin{equation*}
        \text{Fidelity}_{\text{CNOT}} < \text{Fidelity}_{\text{Hadamard}}
    \end{equation*}
    This is one reason why transpilers try to minimize the number of CNOTs and other multi-qubit operations. However, we will understand this in more detail in the next chapters.


    \item \definition{Decoherence}. A \hl{qubit is never perfectly isolated} from the surrounding environment. Even if we do not intentionally measure it, the environment continuously interacts with it: electromagnetic radiation, thermal fluctuations, vibrations, surrounding particles, etc. These \textbf{interactions gradually destroy the quantum properties of the qubit}. In summary, decoherence is the process by which a quantum system progressively loses its ``\emph{quantumness}'' and starts behaving more like a classical object (e.g., a bit instead of a qubit).

    For example, imagine preparing a qubit in the state $\ket{+}$:
    \begin{equation*}
        \ket{+} = \frac{1}{\sqrt{2}}(\ket{0} + \ket{1})
    \end{equation*}
    Initially the qubit exhibits quantum superposition, interference, and phase relationships. However, as time passes, interactions with the environment disturb the phase information. Eventually, the qubit behaves more like a classical probabilistic bit than a true quantum state. The information is not necessarily destroyed instantly, but the quantum coherence gradually leaks into the environment, making it harder to maintain the desired quantum state over time.

    \highspace
    \textcolor{Red2}{\faIcon{exclamation-triangle} \textbf{Why is Decoherence a serious problem?}} A quantum algorithm often requires many sequential operations:
    \begin{equation*}
        U_n U_{n-1} \cdots U_2 U_1
    \end{equation*}
    Where each operation takes a certain amount of time to execute. If the algorithm runs for too long, decoherence may occur before the computation finishes. \textbf{When decoherence becomes dominant}, the superposition is lost, the interference is lost, the entanglement is lost, and the \textbf{quantum advantage disappears}.
\end{itemize}

\highspace
\begin{flushleft}
    \textcolor{Red2}{\faIcon{microchip} \textbf{Hardware Limitations}}
\end{flushleft}
Both gate infidelity and decoherence originate from the same fundamental fact: \textbf{quantum computers are physical devices}. Unlike an ideal mathematical model, qubits are not perfectly isolated, gates are not perfectly accurate, and measurements are not perfectly reliable. Therefore, every real quantum computer has a finite error rate, coherence time, and computational depth. These limitations ultimately restrict how deep a quantum circuit can be before errors become overwhelming.

\highspace
In summary, quantum computers are noisy because the \textbf{gates are imperfectly implemented} and the \textbf{qubits decohere over time}.